\section{Discussion}
\label{sec:discussion}

\subsection{Implications for practice}
\textbf{Stale thresholds are a silent failure.} The most consequential result is not the ranking of the methods
but the UNSW-NB15 failure of the replay baselines (\cref{sec:res-main,sec:res-ablation}). These methods learn the
right classes from the right labels, yet every later flow of those classes is rejected, because the rejection
thresholds were fitted on older traffic. The failure is invisible in the classifier's training loss and in its
detection recall, which stays at 1.0. Any open-world detector whose thresholds are calibrated once is exposed to
it. The remedy is cheap: keep a small store of held-out analyst labels per class and recalibrate on it whenever
the model changes.

\textbf{False positives depend on calibration as much as on the model.} Most benign flows flagged by any method are
\unk{} decisions, not attack labels. Their number is set largely by the rejection quantiles, which all methods
share. A 2.8\% review rate is still large in absolute terms: at one million benign flows a day it means about
28{,}000 reviews, an instance of the base-rate fallacy~\cite{arp2022dos}. \Cref{sec:res-sens} shows how the
quantiles trade review load against recall. Operators should set them from their review capacity, not from a
benchmark.

\textbf{Validate every update.} The deployment guard is the component whose removal costs the most on the
development dataset. A retraining loop that deploys every candidate learns more attack traffic, but it also
deploys models that confuse classes. A hold-out check on the newest labels and a forgetting check on a fixed
reference set are simple to add to any adaptation loop.

\textbf{Read open-world scores with per-class results.} Open-world macro-F1 makes methods that learn different
classes comparable, but it credits a class that is never learned and always rejected as fully correct
(\cref{sec:metrics,sec:res-budget}). A method can therefore win by not learning. We recommend reporting
per-class results and binary detection metrics alongside it, as we do.

\subsection{Threats to validity}
\textbf{Dataset errors.} CIC-IDS2017 and CSE-CIC-IDS2018 contain labelling and feature-extraction
errors~\cite{engelen2021troubleshooting,lanvin2022errors,liu2022error}. Some port-scan flows are feature-identical to
benign flows, and Brute Force and XSS are barely separable. These errors cap every method alike and probably explain
part of the PortScan and web-attack results. Several CSE-CIC-IDS2018 day files are truncated in the published CSVs.
Results on the corrected CIC-IDS2017 release are future work.

\textbf{Short time spans.} The datasets cover two to ten days, so the streams mix attack onsets with modest benign
drift; long-term drift of benign traffic is under-represented. UNSW-NB15 needs held-out classes to contain new
attacks at all.

\textbf{Re-implemented baselines.} The baselines are re-implemented in one multi-class framework with the same
classifier and labelling cost, not run from the authors' code. INSOMNIA and the active-learning baselines were
designed for binary detection or one-shot adaptation, and NetGuard's generative augmentation is not reproduced. Our
results compare the strategies under our protocol; they do not reproduce the original papers' numbers.

\textbf{Simulated analyst.} The oracle returns correct labels with a fixed delay. Real analysts make mistakes,
answer at variable speed, and prefer some flows over others. Label noise would hurt label propagation and the
guard in particular.

\textbf{Design on one dataset.} All design decisions were made on CIC-IDS2017. We checked the class composition of
the other two datasets once, with the static model, which determined the UNSW-NB15 hold-out. The held-out ablations
and the low-budget runs were added after the main results; they are analyses, not design steps.

\textbf{Statistics.} We use three seeds per configuration. The paired Wilcoxon tests treat windows as samples,
but consecutive windows are correlated, so the $p$-values overstate certainty. The seed changes sampling and
initialisation but not the data, so the reported deviations do not include data variability.
